\chapter{Теоретические сведения}
\label{ch:theory}

\section{Электронная почта. Протоколы и архитектура}

\textbf{Электронная почта} (e-mail) --- система обмена текстовыми и мультимедийными
сообщениями через компьютерные сети. Взаимодействие компонентов строится
на трёх ключевых протоколах:
\begin{itemize}
  \item \textbf{SMTP} (Simple Mail Transfer Protocol, RFC~5321) ---
        протокол отправки и пересылки сообщений; порт 25 (сервер--сервер),
        587 (клиент--сервер с аутентификацией);
  \item \textbf{IMAP} (Internet Message Access Protocol, RFC~3501) ---
        протокол доступа к почте с хранением сообщений на сервере;
        порт 143 (STARTTLS), 993 (SSL/TLS);
  \item \textbf{POP3} (Post Office Protocol~3, RFC~1939) ---
        протокол скачивания писем на клиент; порт 110, 995 (SSL).
\end{itemize}

Архитектура почтовой системы включает:
\begin{itemize}
  \item \textbf{MTA} (Mail Transfer Agent) --- агент пересылки; пример: \textit{Postfix};
  \item \textbf{MDA} (Mail Delivery Agent) --- агент доставки в почтовый ящик;
        пример: \textit{Dovecot};
  \item \textbf{MUA} (Mail User Agent) --- клиент пользователя;
        пример: \textit{Roundcubemail}, Thunderbird.
\end{itemize}

\section{iRedMail}

\textbf{iRedMail} --- открытое решение для быстрого развёртывания
полнофункционального почтового сервера. Одним скриптом устанавливаются
и настраиваются следующие компоненты:
\begin{itemize}
  \item \textbf{Postfix} --- MTA (отправка и приём почты);
  \item \textbf{Dovecot} --- IMAP/POP3-сервер и MDA;
  \item \textbf{OpenLDAP} или \textbf{MariaDB}/\textbf{PostgreSQL} ---
        хранилище учётных записей;
  \item \textbf{Roundcubemail} --- веб-интерфейс почты;
  \item \textbf{iRedAdmin} --- панель администрирования;
  \item \textbf{Amavisd-new}, \textbf{SpamAssassin}, \textbf{ClamAV} ---
        антиспам и антивирус;
  \item \textbf{Fail2ban} --- защита от брутфорса;
  \item \textbf{Nginx} или Apache --- веб-сервер.
\end{itemize}

Версия, использованная в работе: \textbf{1.6.2}.

\section{OpenLDAP как бэкенд}

\textbf{LDAP} (Lightweight Directory Access Protocol) --- протокол доступа
к каталогам. iRedMail использует OpenLDAP для хранения учётных записей
пользователей и доменов. Суффикс (root DN) дерева каталога определяется
по имени почтового домена, например:
\begin{lstlisting}
dc=yazikov,dc=iks531,dc=local
\end{lstlisting}

\section{Требования к хосту}

iRedMail предъявляет строгие требования к имени хоста:
\begin{itemize}
  \item FQDN должен разрешаться в IP-адрес сервера (A-запись в DNS);
  \item FQDN \textbf{не должен} совпадать с именем почтового домена:
        если домен \texttt{yazikov.iks531.local}, то FQDN должен быть
        \texttt{mail.yazikov.iks531.local};
  \item \texttt{/etc/hosts} должен содержать строку с IP и FQDN;
  \item \texttt{hostname -f} должен возвращать полное доменное имя.
\end{itemize}

\section{SSL/TLS-сертификат}

iRedMail автоматически генерирует \textbf{самоподписанный SSL-сертификат}
для шифрования SMTP, IMAP и HTTPS-трафика. Браузер при первом открытии
выдаст предупреждение о недоверенном сертификате --- это ожидаемо
для учебного стенда. Для принятия необходимо нажать
\textit{Advanced} → \textit{Accept the Risk and Continue}.

\endinput
